% One-time imported draft. Editable native source; no Markdown regeneration.
\section{Step 13 Compare the full register
  values}\label{004-bit-inspector-s13-step-13-compare-the-full-register-values}

The four formatting services report a word, so they cannot by themselves
reveal a difference in the high register half. The full register view
and the stored doublewords supply that evidence.

\subsection{Write or
  inspect}\label{004-bit-inspector-s13-write-or-inspect}

Debug the default input and stop at values\_stored. Compare
\texttt{\$s0} with \texttt{\$s1}, then compare the corresponding
eight-byte fields. Resume and check the full output.

\subsection{Expected
  observation}\label{004-bit-inspector-s13-expected-observation}

\texttt{\$s0} and \texttt{\$s1} have different high halves, while both
have low word ffffffff. Passing either as the word-format argument gives
the same report.

\subsection{Check your
  understanding}\label{004-bit-inspector-s13-check-your-understanding}

Does matching hexadecimal console output prove two 64-bit registers are
equal?

\begin{itemize}
  \tightlist
  \item
        \textbf{A} Yes; eight digits always describe a 64-bit register.
  \item
        \textbf{B} Yes; the service automatically prints hidden upper digits.
  \item
        \textbf{C} No; the service only reports their low words.
\end{itemize}
